\chapter{Standard Model neutrinos}
\label{chapter:nuIntro}
Neutrinos are elementary particles, specifically neutral leptons, that, according to the \ac{SM} of Particle Physics \cite{fundamentalsOFnu}, only interact via weak interactions.  
They were first postulated in 1930 by the Austrian physicist Wolfgang Pauli to describe the continuous energy spectrum of beta decay. 
Since only the electron and the recoiling nucleus were observed, the process seemed to violate energy and angular momentum conservation; therefore, Pauli introduced a third, 
unseen spin-1/2 particle to carry away the missing energy and balance the angular momentum \cite{Pauli:1930neutrinos}.

They were eventually confirmed experimentally in 1956 by Clyde Cowan and Frederick Reines during an experiment conducted at the Savannah River 
nuclear reactor in South Carolina. 
This experiment exploited inverse beta decay, in which the postulated particle---specifically an electron antineutrino ($\overline{\nu}_e$)---interacts with a proton to produce a neutron and a positron, as described by Eq.~\ref{eq:inverse_beta_decay}.
\begin{equation}
    \label{eq:inverse_beta_decay}
    \overline{\nu}_e + p \rightarrow n + e^{+}
\end{equation}
In the experimental setup, antineutrinos produced by the reactor fission processes interacted with protons in the water of two target tanks, 
yielding a positron and a neutron. The positron quickly annihilated with an adjacent electron, producing two prompt gamma rays. 
Meanwhile, the neutron thermalized and was subsequently captured by the cadmium nuclei of cadmium chloride (CdCl$_2$) dissolved in the water tanks.
As a consequence of this neutron capture, an additional gamma-ray signal was released. The detection of these three gamma-ray signals 
in delayed coincidence provided an unambiguous signature of the antineutrino interaction, definitively confirming the existence of 
both neutrinos and antineutrinos~\cite{cowan1956nudetection}.

Within the framework of the \ac{SM}, neutrinos are described as neutral fermions that do not participate in strong or 
electromagnetic interactions: they are colour singlets and carry no electric charge. The \ac{LH} neutrino of each family is grouped with the corresponding charged lepton ($\ell = e, \mu, \tau$) into a weak-isospin doublet $L_\ell = (\nu_\ell, \ell)^T_\text{L}$.
% , which transforms as $(\mathbf{1}, \mathbf{2}, -\tfrac{1}{2})$ under the gauge group SU$(3)_\text{C} \times$ SU$(2)_\text{L} \times$ U$(1)_Y$. The electromagnetic group U$(1)_\text{EM}$ is not part of the original gauge group: it is the subgroup that survives electroweak symmetry breaking, with $Q = T_3 + Y$. 
No \ac{RH} neutrino field is present in the minimal \ac{SM}.

The SU$(2)_\text{L}$ gauge invariance fixes the form of the weak interactions between neutrinos and the gauge bosons; neutrinos taking part in these interactions are called \textit{active} and, in the minimal \ac{SM}, are massless. These interactions consist of \ac{CC} processes mediated by the $W^\pm$ bosons and \ac{NC} processes mediated by the $Z^0$ boson, described by the following Lagrangians:
\begin{equation}
    \label{eq:SM-neutrino-lagrangian}
    \begin{aligned}
        \mathcal{L}_\text{CC} &= - \frac{g}{\sqrt{2}} \sum_{\ell} \overline{\nu}_{\text{L},\ell} \gamma^\mu \ell_\text{L}^- W_\mu^+ + \text{h.c.} \\
        \mathcal{L}_\text{NC} &= - \frac{g}{2 \cos \theta_W} \sum_{\ell} \overline{\nu}_{\text{L},\ell} \gamma^\mu \nu_{\text{L},\ell} Z_\mu^0
    \end{aligned}
\end{equation}
where $g$ is the SU$(2)_\text{L}$ coupling constant and $\theta_W$ is the Weinberg weak mixing angle. 

Precision measurements of the invisible decay width of the $Z^0$ boson at LEP constrain the number of active neutrino families to $N_\nu = 2.9963 \pm 0.0074$~\cite{JanotJadach2020}, consistent with the three known generations of charged leptons\footnote{The measurement only counts active neutrinos lighter than $m_Z/2$; it does not exclude sterile neutrinos.}.